\chapter{Regulation and expression across levels}
Two cells can carry the same gene and encounter the same present signal,
yet contain different amounts of its protein. A model that chooses one
stored function from the current input cannot explain that difference unless
it also represents history. This is the first useful boundary between
biological regulation and Chapter 10's instantaneous expert routing.

We will connect a promoter-state model to RNA and protein dynamics,
derive a reproducible update, and ask which parameters observations can
identify. Then we will extract a software contract without claiming that a
neural module is a molecule. All numbers are assigned teaching values,
not measurements or an Evolutor biological-performance result.

\section{Stored sequence is not executed function}
\index{regulation}\index{expression}\index{promoter}
DNA supplies a substrate and information within cellular machinery.
Transcription requires machinery interacting with DNA; translation requires
RNA, ribosomes and other resources. Existing RNA and protein can persist
after production changes. Protein abundance is distinct from protein
activity, which may depend on localization, partners or chemical state.
These distinctions make \emph{turning a gene on} ambiguous: does it mean
promoter occupancy, RNA synthesis, RNA abundance, protein synthesis,
protein abundance, or a downstream effect?

Our setting is a simplified bacterial promoter with one overlapping repressor
site. Simultaneous polymerase/repressor occupancy is excluded.
Garcia and Phillips analyze such a three-state simple-repression motif under
quasi-equilibrium and promoter-occupancy assumptions~\cite{garcia}.
We use that limited starting point, not their fitted parameters or experimental
performance.
\Plate{01-promoter}{Empty, polymerase-bound and repressor-bound promoter
states. The overlapping binding region explains exclusion in this model.
DNA remains present in all states; only polymerase-bound occupancy is assigned
transcriptional activity. Molecular silhouettes are schematic.}{fig:promoter}

\section{Normalize states before interpreting a gate}
\index{statistical weight}
Assign dimensionless weights $1,u,r$ to the three states. These summarize
effective availability and binding preference relative to an empty promoter.
They are not probabilities until normalized:
\begin{equation}
 Z=1+u+r,\qquad P_{\emptyset}=Z^{-1},\quad
 P_{\mathrm{pol}}=u/Z,\quad P_{\mathrm{rep}}=r/Z.
\end{equation}
If transcription is proportional to polymerase occupancy, the ratio with
and without repressor is, for $u>0$,
\begin{equation}
 f(r)=\frac{u/(1+u+r)}{u/(1+u)}=\frac{1+u}{1+u+r}.
\end{equation}
Only for $u\ll1$ is $f(r)\approx1/(1+r)$ a suitable approximation.
At $u=0$ both transcription rates vanish: the measured ratio is undefined,
even though the simplified expression has a finite limit.
\Listing{occupancy}
\begin{center}\input{\Generated/occupancy-table.tex}\end{center}
At $u=0.2,r=1$, active probability is $1/11$ and fold change is $6/11$.
Confusing these quantities changes the interpretation of an experiment.
The code scales weights before summation to avoid overflow when two finite
weights are very large.

Unlike a generic softmax, the alternatives here are specified physical
configurations. But many possible states and processes have been omitted.
A cooperative mechanism needs its allowed configurations and weights;
raising $r$ to an arbitrary power does not itself derive cooperativity.
A fitted response curve need not uniquely identify a mechanism.

\section{Track production, retention and removal}
\index{transcription}\index{translation}\index{turnover}
Let $m(t),p(t)$ be coarse-grained RNA and protein concentrations in assigned
nM units. Time is minutes, RNA production $s(t)$ is nM/min,
translation $k$ is $\mathrm{min}^{-1}$, and effective loss rates
$\gamma_m,\gamma_p>0$ are $\mathrm{min}^{-1}$. With fixed volume and
first-order removal, define
\begin{align}
 \dot m&=s(t)-\gamma_m m,\\
 \dot p&=km-\gamma_p p.
\end{align}
Translation does not consume RNA in this model. Each RNA can support repeated
production; degradation and dilution enter separately through effective loss.
This mean-concentration model omits discrete fluctuations and does not
describe each individual cell.
\Plate{02-expression}{The objects behind the two-stage model. Polymerase
uses DNA to synthesize RNA; a ribosome uses RNA while a peptide grows.
Loss removes RNA and protein from their respective pools. This coarse
bacterial mechanism omits detailed resources and regulatory factors.}{fig:expression}

Introduce $\alpha$ in nM/min and a competent-promoter fraction
$a(t)\in[0,1]$, then set
\begin{equation}
 s(t)=\alpha a(t)\frac{u(t)}{1+u(t)+r(t)}.
\end{equation}
Occupancy is conditioned on the competent subpopulation. The product treats
$a$ as separately supplied; it is not a derivation of chromatin accessibility
or a universal bacterial rule. Gene copy number and fixed-volume conversion
are absorbed into $\alpha$. Altering them changes the model.

\section{Solve a constant-input interval}
\index{exact propagation}
Hold all coefficients fixed for duration $\Delta$, and set $m_*=s/\gamma_m$.
Multiplying the RNA equation by $e^{\gamma_m t}$ and integrating yields
\begin{equation}
 m(\Delta)=m_*+(m_0-m_*)e^{-\gamma_m\Delta}.
\end{equation}
For protein, use the integrating factor $e^{\gamma_p t}$:
\begin{align}
 p(\Delta)&=p_0e^{-\gamma_p\Delta}
  +k\int_0^\Delta e^{-\gamma_p(\Delta-\tau)}m(\tau)\,d\tau\\
 &=p_0e^{-\gamma_p\Delta}
  +k\left[m_*\frac{1-e^{-\gamma_p\Delta}}{\gamma_p}
    +(m_0-m_*)J(\Delta)\right],
\end{align}
where
\begin{equation}
 J(\Delta)=\begin{cases}
 \displaystyle\frac{e^{-\gamma_m\Delta}-e^{-\gamma_p\Delta}}
 {\gamma_p-\gamma_m},&\gamma_p\ne\gamma_m,\\[6pt]
 \Delta e^{-\gamma_m\Delta},&\gamma_p=\gamma_m.
 \end{cases}
\end{equation}
At equal rates, the integrand's $\tau$ dependence cancels. The apparent
singularity is removable, not a physical divergence.

For $g=|\gamma_p-\gamma_m|>0$, rewrite
$J=e^{-\min(\gamma_m,\gamma_p)\Delta}(1-e^{-g\Delta})/g$.
This avoids subtracting nearly equal exponentials and avoids a large
positive exponential. The implementation uses \texttt{expm1} for
$1-e^{-x}$. An exact mathematical formula still needs floating-point tests.
\Listing{expression_step}
The immutable \texttt{State} stores RNA/protein concentrations; immutable
\texttt{Definition} stores four rate parameters. Both live in
\texttt{reference.py}. The context wrapper is:
\Listing{regulated_step}
There is no winner selection or capacity budget here. The wrapper updates
a continuous state. Multiple intervals equal one long interval only with
unchanged input and parameters, additive durations, and carried state.
Tests verify this semigroup property and compare with an independent
small-step fourth-order Runge--Kutta integration.

\section{A stopped source can leave a rising output}
\index{memory}\index{step response}
Take $\alpha=12,u=0.2,r=0,a=1,k=2,\gamma_m=1,\gamma_p=0.5$,
so $s=2$. Start at zero and set $a=0$ after four minutes, retaining
the RNA/protein state. The reference uses exact quarter-minute updates:
\Listing{simulate}
\Plate{03-response}{RNA and protein under an assigned four-minute production
pulse. RNA falls after closure, but protein briefly continues rising because
surviving RNA still drives translation. Solid and dashed lines encode the
two quantities independently of color.}{fig:response}
At four minutes, $m\approx1.96337,p\approx5.98116$. Immediately after
closure, $\dot m=-1.96337$ nM/min, yet
$\dot p=2(1.96337)-0.5(5.98116)\approx0.93616$ nM/min.
Rising protein does not establish current transcriptional activity.

Setting derivatives to zero for constant source gives
$m_*=s/\gamma_m$ and $p_*=ks/(\gamma_m\gamma_p)$:
\Listing{steady_state}
The assigned steady state is $(2,8)$ nM. Closing its source for one minute
gives approximately $(0.73576,6.76146)$, whereas zero initial state under
the identical closed context stays zero. Context alone cannot determine
abundance. Production history is part of the computational state.

\section{A perfect output trajectory can remain ambiguous}
\index{identifiability}\index{intervention}
Consider three parameter sets, all with $\gamma_p=0.5$:
A has $(s,k,\gamma_m)=(2,2,1)$; B has $(1,4,1)$; C has $(4,2,2)$.
All have $p_*=8$. A and C approach it differently, so a transient can
distinguish them. But A and B have identical protein trajectories from
zero despite different RNA trajectories.
\begin{center}\input{\Generated/identifiability-table.tex}\end{center}
\Plate{04-identifiability}{A and C share a steady protein value but differ
during the transient. A and B coincide throughout: their transcription--translation
product and loss rates match. RNA observations can separate A from B.}{fig:identifiability}
For any $c>0$, replace $s$ by $s/c$, $k$ by $ck$, and $m_0$ by $m_0/c$.
The RNA solution is divided by $c$, preserving $km(t)$.
With unchanged $p_0$, the entire protein solution is unchanged.
More precise protein measurements cannot remove this structural ambiguity.
RNA measurements, an independent translation constraint, or an intervention
that breaks the symmetry are needed.

Low protein therefore does not uniquely establish transcriptional repression.
RNA turnover, translation, protein turnover and history offer alternatives.
A genomic predictor should name its output and the evidence that distinguishes
these explanations. Predictive association and mechanistic intervention are
different tests.

\section{Extract a typed software hypothesis}
\index{state ownership}\index{biological analogy}
Separate a stored definition, context, eligibility, evolving state and executed
effect. A definition is not a state; eligibility need not guarantee admission
under a compute budget; abundance need not equal activity; an output is not
the specification. One variable called a gene cannot safely stand for all five.
\Plate{05-contract}{Equal present context with different initial state
produces different next state. Solid arrows are update dependencies.
Dashed boundaries mark state owned by one simulated system, not global
state shared across unrelated records.}{fig:contract}

Chapter 10's gate weights a current-token expert branch under cohort and
capacity rules. This chapter retains material/state after the source closes.
Naming both mechanisms regulation does not make their semantics equivalent.
A regulation-inspired software experiment must specify update time, reset
boundaries, parameter ownership, observations and interventions. Accidental
state carryover between unrelated genomic records is leakage, not useful memory.

These are research hypotheses above either Evolutor direction: Hermon DNA
remains Transformer-based and DOGMA remains non-Transformer.
The scalar model establishes no biological capability for either.
A useful next study would compare stateful and stateless predictors with
matched data and compute budgets, held-out interventions and explicit
reset ablations. Chapter 12 asks a different question: how can a compact
specification generate a larger computational structure before execution?

\clearpage
\section{Exercises with worked reasoning}
\begin{enumerate}
\item \textbf{Normalization.} Weights $(1,.2,1)$ give which active probability?
 \emph{Solution:} $.2/2.2=1/11$.
\item \textbf{Fold change.} Is $1/11$ the fold change?
 \emph{Solution:} no; divide by unrepressed $1/6$ to obtain $6/11$.
\item \textbf{Zero polymerase.} Interpret $u=0$.
 \emph{Solution:} no active occupancy; the ratio of zero rates is undefined.
\item \textbf{Units.} Check $km$.
 \emph{Solution:} $\mathrm{min}^{-1}\cdot\mathrm{nM}=\mathrm{nM/min}$.
\item \textbf{Equal rates.} With both loss rates one, $s=2,k=3$ and zero
 state, derive $p(t)$. \emph{Solution:} $6[1-(1+t)e^{-t}]$.
\item \textbf{Steady state.} Find the pulse's active-source equilibrium.
 \emph{Solution:} $(2,8)$ nM.
\item \textbf{Delayed peak.} Why can protein rise after source closure?
 \emph{Solution:} $km$ can still exceed $\gamma_p p$.
\item \textbf{Memory.} Compare cold and warm states one minute after closure.
 \emph{Solution:} $(0,0)$ versus approximately $(.73576,6.76146)$.
\item \textbf{Ambiguity.} Can protein-only time series distinguish A and B?
 \emph{Solution:} no; the scaling symmetry preserves all those observations.
\item \textbf{Chunking.} When does two-step propagation equal one step?
 \emph{Solution:} same held coefficients, additive durations and carried state.
\item \textbf{Independent check.} Why also compare against RK4?
 \emph{Solution:} the printed formula and its implementation might share a mistake.
\item \textbf{Study design.} Test a regulation-inspired predictor.
 \emph{Rubric:} name quantities, reset boundaries, matched baselines,
 intervention splits, uncertainty and competing explanations.
\end{enumerate}
\section*{Selective glossary}
\addcontentsline{toc}{section}{Selective glossary}
\textbf{Occupancy:} probability of a molecular configuration.
\textbf{Expression:} a process whose measured level must be specified.
\textbf{Turnover:} effective removal from a modeled pool.
\textbf{Identifiability:} distinct parameters imply distinguishable observations.
\textbf{Context:} supplied update conditions, distinct from carried state.
